\subsubsection{Transport électronique bidirectionnel et registre arithmétique}
\label{ii:transporte-electronico}

La représentation centrale utilise une cellule partagée par les deux
évaluations APP. Ses coordonnées appartiennent à \(D_9=\{1,\ldots,9\}\),
et l'inversion \((r,c)\mapsto(10-r,10-c)\) possède l'unique point fixe
\((5,5)\). L'état visible s'écrit \((r,c,h,v,n)\), avec
\(h,v\in\{-1,1\}\) et un indice chronologique entier \(n\geq0\).
Le sélecteur tritique répète \((+1,-1,0)\). Dans le feuillet additif, on lit
\(r+c\) et l'on translate la colonne dans la direction \(h\) ; dans le
feuillet multiplicatif, on lit \(rc\) et l'on translate la ligne dans la direction
\(v\). Le seuil conserve la position. Les deux orientations sont
inversées à l'achèvement de chaque intervalle de vingt-sept mises à jour.

Pour exprimer la lecture et la retenue arithmétique dans la même convention positive,
définissons
\[
 [z]^+_9=1+((z-1)\bmod9),\qquad
 q_9(z)=\frac{z-[z]^+_9}{9}.
\]
L'évaluation précède le déplacement. Si \(a_n\) est l'entier
lu, l'événement conserve \((a_n,[a_n]^+_9,q_9(a_n))\), le feuillet,
la position précédente, l'orientation et le franchissement de frontière. Le
symbole tritique est \([a_n]^+_9\bmod3\). Le relèvement spatial
conserve, par exemple, l'entier
\(q_9(c+h)\) lorsque \(c\) est remplacé par \([c+h]^+_9\).

La représentation à cellule partagée et la représentation à deux
curseurs de \eqref{ii:tpk-fobs-posiciones} possèdent des états visibles
différents. Leurs registres sont reliés à l'état TPK par
leurs applications de représentation respectives. Pour la trajectoire centrale,
on utilise ici intégralement la mise à jour sur la cellule
partagée qui produit le mot électronique.

\begin{proposition}[Mise à jour inversible avec quotients conservés]
\label{ii:electron-inversa-cocientes}
Ajoutons à l'état visible deux registres entiers \(Q_+,Q_\times\)
et deux nombres d'enroulement \(W_r,W_c\). Chaque lecture active
incrémente le registre de son feuillet de \(q_9(a_n)\) ; chaque déplacement
incrémente l'enroulement correspondant de son franchissement de frontière.
Cette mise à jour possède un inverse explicite sur son image.
\end{proposition}
\begin{proof}
L'indice suivant détermine le mode précédent et indique si une inversion
des orientations vient de se produire. Cette inversion est d'abord annulée.
Si le mode était additif, la colonne précédente est retrouvée sous la forme
\([c-h]^+_9\) ; s'il était multiplicatif, la ligne est retrouvée sous la forme
\([r-v]^+_9\). Le mode neutre conserve les deux coordonnées. La cellule
retrouvée détermine de nouveau l'entier lu, son résidu, son
quotient et le franchissement de frontière. En soustrayant ces incréments de
\(Q_+,Q_\times,W_r,W_c\) et en inversant l'indice chronologique, on
retrouve exactement l'état précédent. La composition de ces
inverses permet de retrouver une histoire finie complète. La concaténation des
événements conserve en outre l'ordre des lectures.
\end{proof}

Les registres \(Q_+,Q_\times\) représentent l'historique des quotients au moyen de sommes explicites. La construction s'intègre dans la fibre de mémoire avec ses autres coordonnées. Sa loi de composition, pour des historiques concaténables \(u,v\), est
\[
 Q_s(uv)=Q_s(u)+Q_s(v),\qquad s\in\{+,\times\},
\]
où la seconde histoire commence dans l'état final de la première.
L'inversion algébrique d'une mise à jour soustrait ses incréments ;
parcourir ensuite un trajet spatial de retour enregistre
de nouvelles évaluations. Cette distinction conserve simultanément la
réversibilité du transport et la durée de l'histoire.

\begin{proposition}[Retour central et accumulation exacte]
\label{ii:electron-memoria-retorno}
À partir de \((r,c,h,v,n)=(5,5,1,1,0)\) et de registres initialement nuls, la
projection spatiale et ses orientations reviennent après cinquante-quatre
mises à jour. Dans cet intervalle,
\[
 (\Delta W_r,\Delta W_c)=(0,0),\qquad
 (\Delta Q_+,\Delta Q_\times)=(10,46).
\]
Après \(54N\) mises à jour, pour tout entier \(N\geq0\),
on a \((Q_+,Q_\times)=(10N,46N)\), tandis que les coordonnées
spatiales et leurs orientations reviennent. L'indice entier prend la valeur
\(54N\) ; le calendrier dodécaphasique modulo \(108\) revient pour
\(N\) pair.
\end{proposition}
\begin{proof}
Chaque couple de modes actifs déplace une fois chaque coordonnée. En
vingt-sept mises à jour, les deux parcourent un tour de neuf
positions et reviennent à la diagonale centrale ; \(h,v\) sont ensuite
inversés. Le bloc suivant parcourt le tour opposé. Les franchissements
spatiaux s'annulent entre les deux blocs.

Dans chaque bloc, les lectures additives parcourent \(2k\),
\(1\leq k\leq9\), dont les quotients positifs ont pour somme cinq. Les lectures
multiplicatives du bloc positif parcourent
\(k[k+1]^+_9\). Leur somme entière et la somme de leurs résidus sont
\[
 \sum_{k=1}^9 k[k+1]^+_9=249,\qquad
 \sum_{k=1}^9[k[k+1]^+_9]^+_9=42.
\]
La seconde identité résulte de la liste
\((2,6,3,2,3,6,2,9,9)\). Par division euclidienne,
la somme des quotients est \((249-42)/9=23\). La substitution
\(k\mapsto[k-1]^+_9\) montre que le bloc négatif contient
le même multiensemble de produits. Les deux blocs produisent donc
dix et quarante-six unités dans les registres respectifs.
La projection spatiale et ses orientations reviennent à leurs valeurs
initiales. Les règles d'activation et d'inversion se répètent lorsqu'on
décale l'indice de \(54\), et les incréments ne dépendent pas des
registres accumulés. La loi additive
précédente démontre par récurrence la formule pour tout \(N\).
\end{proof}

La lecture tritique et l'orientation par fenêtres fournissent, sur
la même trajectoire,
\[
 \gamma_e=\bigl((100000220)^3(120200000)^3\bigr)^2,\qquad
 \tau_e=(+,+,+,-,-,-,+,+,+,-,-,-).
\]
Les puissances représentent des concaténations. Les lecteurs d'écarts et de fenêtres développés dans \eqref{eq:el-kd} et \eqref{eq:el-komega} calculent tous deux \((80,54,6)\). L'évaluation s'obtient à partir de la trajectoire avant d'incorporer \(\alpha\), la section d'action ou une masse. La mémoire des quotients conserve des informations supplémentaires par rapport à ce triplet : en cent huit mises à jour, elle enregistre \((20,92)\), avec un enroulement spatial net nul. Chaque représentation maintient ainsi le type de l'observable qui la produit.
